\subsection{Fixed-parameter FlexQP reconstruction}\label{sec:flex}
FlexQP addresses elastic quadratic optimization with a splitting formulation;
Deep FlexQP additionally uses a learned numerical-parameter policy~\cite{flex}.
The implemented reference reconstructs the fixed equations and Algorithm~1 from
the paper. It is neither the authors' released implementation nor a reproduction
of a learned policy. Trained weights, training data and a reproducible policy
implementation were unavailable within the comparison. The fixed reconstruction
is nevertheless executable on its stated domain and is evaluated against the
same independent lifted original objective.

The admitted model has positive finite-side linear prices and no finite hard
rows or quadratic soft costs. After converting finite sides to inequalities,
the elastic problem can be written as
\begin{equation}\label{eq:flex-model}
 \min_{x\in\R^n,\ \xi\geq0}\ \tfrac12x^\top P x+q^\top x+\mu^\top\xi
 \quad\text{subject to}\quad Gx-\xi\leq h,\qquad \mu>0.
\end{equation}
Lower interval sides become signed inequalities with their own prices. This
conversion preserves directionality, but it does not justify giving a hard
physical row a finite weight. The adapter consequently refuses hard rows rather
than choosing a numerically large weight and calling it exact hardness.

The fixed numerical parameters are $\sigma_x=10^{-6}$, $\sigma_s=1$, $\rho=1$
and relaxation $\alpha=1.6$. Eliminating the slack-related block produces the
Schur system
\begin{equation}\label{eq:flex-system}
 K=P+\sigma_x I+cG^\top G,\qquad
 c=(\sigma_s^{-1}+\rho^{-1})^{-1}.
\end{equation}
The positive proximal terms stabilize the numerical subproblem. They do not
modify the scientific objective in \eqref{eq:flex-model}; treating $\sigma_x$
as an objective regularizer would change the target optimizer. This distinction
is checked through original objective evaluation and stationarity, not through
a comparison of the stabilized subproblem's objective.

The reconstruction uses \texttt{softqp.flex.FlexWorkspace} in SoftQP 0.1.0
at the native source revision in the reproducibility section, with SciPy
1.18.1 for the dense Cholesky factor and solves. The experiment uses the workspace's absolute and relative tolerances of $10^{-7}$
and an iteration limit of 200000. The proximal parameters above remain
fixed across the cases. For $k$ finite sides and $n$ decisions, this dense
implementation uses $O(kn+n^2)$ memory, $O(kn^2+n^3)$ preparation and
$O(kn+n^2)$ arithmetic per iteration.
It is less mature and has different language-level overhead from the compiled
C backends. Observed runtime differences therefore measure these concrete
implementations under the protocol, rather than an implementation-independent
lower bound on the FlexQP formulation. The reconstruction is not retuned after
examining the held-out results. In particular, the comparison does not choose a
different proximal parameter for each unfavorable test instance.

The split uses a nonnegative margin $s$ and a signed elastic coordinate
$z$, with $Gx+s-z=h$ and price $\sum_i\mu_i|z_i|$. For a fixed $x$, each
row chooses $z_i=(G_i x-h_i)_+$ and $s_i=(h_i-G_i x)_+$ at minimum
price, recovering \eqref{eq:flex-model}. The signed variable $z$ in this
split is distinct from an original interval-row value.
Let $w$ and $y$ be the split multipliers for the margin and elastic
coordinates. Starting from zero for a cold call, a step computes
\begin{align}\label{eq:flex-graph}
 K\widetilde x&=\sigma_x x-q+
   cG^\top(h-s+w/\sigma_s+z-y/\rho),\\
 \nu&=c(G\widetilde x+s-w/\sigma_s-z+y/\rho-h),\nonumber\\
 \widetilde s&=s-w/\sigma_s-\nu/\sigma_s,\qquad
 \widetilde z=z-y/\rho+\nu/\rho.\nonumber
\end{align}
For $d=x,s,z$ put $\bar d=\alpha\widetilde d+(1-\alpha)d$.
Writing $\mathcal T_{\tau}(v)_i=\operatorname{sign}(v_i)
(|v_i|-\tau_i)_+$ for componentwise soft thresholding, the remaining
updates are
\begin{align}\label{eq:flex-project}
 x^+&=\bar x,\qquad s^+=(\bar s+w/\sigma_s)_+,\\
 z^+&=\mathcal T_{\mu/\rho}(\bar z+y/\rho),\nonumber\\
 w^+&=w+\sigma_s(\bar s-s^+),\qquad
 y^+=y+\rho(\bar z-z^+).\nonumber
\end{align}
Here the superscript $+$ denotes the next iterate, while a subscript $+$
denotes a positive part. The solver checks the original-space residuals
\begin{equation}\label{eq:flex-stop}
 r_{\rm d}=\norm{Px+q+G^\top y},\qquad
 r_{\rm p}=\frac{\norm{y-\Pi_{[0,\mu]}(y+\rho(Gx-h))}}{\rho}.
\end{equation}
The dual tolerance scales its relative part by
$\max(\norm{Px},\norm{q},\norm{G^\top y})$. The projected-box tolerance
uses $\max(\norm{Gx},\norm{h})$. Both add the absolute tolerance.
This checks stationarity and capped inequality complementarity rather
than just the size of an iterate change. It does not imply convergence
for arbitrary adaptive parameters or a finite-precision rate.

\begin{algorithm}[htbp]
\caption{Execute the fixed FlexQP equations on an admitted pure linear elastic problem, preserving original prices and independently testing the returned original-space optimum.}\label{alg:flex}
\begin{algorithmic}[1]
\Require Compatible pure linear elastic inequalities and fixed numerical parameters
\State Refuse finite hard rows, quadratic prices and nonpositive admitted side prices
\State Convert finite interval sides into $Gx-\xi\leq h$ with original linear prices
\State Prepare \eqref{eq:flex-system} using the fixed proximal parameters
\State Set the cold state $(x,s,z,w,y)$ to zero
\While{the iteration limit is not reached}
\State Solve \eqref{eq:flex-graph}, then relax and apply \eqref{eq:flex-project}
\State Refuse a nonfinite decision and test \eqref{eq:flex-stop} against both tolerances
\State Stop on convergence and retain a distinct status at the iteration limit
\EndWhile
\State Recover interval dual signs and compare the original objective to lifting
\end{algorithmic}
\end{algorithm}

The relationship between a finite elastic optimum and hard feasibility requires
care. If the original hard inequalities are feasible and suitable multiplier
bounds are known, a sufficiently large exact linear price can recover a hard
solution. That result does not make every finite price exact, nor does it make
the output safe for a plant whose immutable limits have been softened. The
application explicitly describes the objective being priced and separates it
from the original hard mask. This is a modeling and acceptance requirement,
not a criticism unique to one elastic algorithm.

Compatibility exclusions remain part of the scientific denominator. Symmetric
and asymmetric elastic families with nonzero quadratic side prices are outside
this pure linear reconstruction. They are marked incompatible before timing,
not counted as inaccurate solves or assigned zero durations. Conversely,
admitted asymmetric linear cases are executed and checked, so an
unavailable learned policy is not used to excuse omission of the fixed published
reference. The complete six-path original comparison can therefore expose
both semantic limits and implementation costs without pretending that all
methods solve every penalty family.
The factor in the Schur solve depends on the Hessian, signed inequality
matrix and proximal settings. A vector-only change can leave that factor
invariant mathematically. The implemented \path{FlexWorkspace} deep-copies
one problem at initialization and exposes reset and solve for that fixed
problem. Changing the vectors requires a new workspace. The dense
reconstruction caps admission at 2048 decision variables and 8192
finite sides. The complete-call timer
includes admission, setup and original-space recovery. The theoretical proximal regularizers
and the practical implementation overhead must both remain visible when
interpreting a comparison with a mature compiled backend.
